\documentclass[11pt]{article}

\usepackage[a4paper,margin=1in]{geometry}
\usepackage{amsmath,amssymb,amsthm,mathtools}
\usepackage{booktabs}
\usepackage{tabularx}
\usepackage{enumitem}
\usepackage{microtype}
\usepackage[hidelinks]{hyperref}
\usepackage{xcolor}

\newtheorem{theorem}{Theorem}[section]
\newtheorem{proposition}[theorem]{Proposition}
\newtheorem{lemma}[theorem]{Lemma}
\newtheorem{corollary}[theorem]{Corollary}
\theoremstyle{definition}
\newtheorem{definition}[theorem]{Definition}
\newtheorem{problem}[theorem]{Problem}
\newtheorem{remark}[theorem]{Remark}

\newcommand{\psd}{\mathbb{H}_{+}}
\newcommand{\Tr}{\operatorname{tr}}
\newcommand{\diag}{\operatorname{diag}}
\newcommand{\smin}{s_{2}}
\newcommand{\smax}{s_{1}}
\newcommand{\lamax}{\lambda_{1}}
\newcommand{\lamin}{\lambda_{2}}
\newcommand{\statusL}{\textcolor{blue!70!black}{[L]}}
\newcommand{\statusR}{\textcolor{green!50!black}{[R]}}
\newcommand{\statusW}{\textcolor{orange!80!black}{[W]}}
\newcommand{\statusN}{\textcolor{purple!70!black}{[N]}}
\newcommand{\statusF}{\textcolor{red!70!black}{[F]}}

\title{The Three-Factor Matrix Young Problem in Dimension Two\\
\large Working manuscript, rigorous reductions, and proof audit}
\author{Working research notes}
\date{7 August 2026}

\begin{document}
\maketitle

\begin{abstract}
We study the unresolved low-dimensional cell of the multivariable Matrix Young conjecture with three factors and matrix size two.  Pang recently showed that the conjecture fails for every $m\ge4$ already for $2\times2$ real rank-one projections, and also fails for $m=3$ in dimension $3$, at the second singular-value/eigenvalue index.  Thus $(m,n)=(3,2)$ is a natural boundary case.  These notes collect the work completed so far, with a strict status audit.  In particular, we separate literature results, internally rigorous reductions, working lemmas that still require an independent proof audit, numerical evidence, and routes that have already been falsified.
\end{abstract}

\paragraph{Status convention.}
\statusL{} = established literature result; \statusR{} = rigorous derivation retained in the present notes; \statusW{} = working lemma or proof chain not yet independently closed; \statusN{} = numerical/experimental evidence only; \statusF{} = route falsified or withdrawn.

\tableofcontents

\section{Problem and literature boundary}

Let $A_1,A_2,A_3\in\mathbb H_+^n$ and let $p_i>1$ satisfy
\[
\frac1{p_1}+\frac1{p_2}+\frac1{p_3}=1.
\]
Write
\[
\alpha=\frac1{p_1},\qquad \beta=\frac1{p_2},\qquad \gamma=\frac1{p_3},
\qquad \alpha+\beta+\gamma=1,
\]
and set
\[
X=A_1^{p_1},\qquad Y=A_2^{p_2},\qquad Z=A_3^{p_3}.
\]
The three-factor Matrix Young problem asks whether
\begin{equation}
 s_j\!\left(X^\alpha Y^\beta Z^\gamma\right)
 \le
 \lambda_j\!\left(\alpha X+\beta Y+\gamma Z\right),
 \qquad j=1,\dots,n.
 \label{eq:MY}
\end{equation}
For $m=2$, Ando's matrix Young inequality gives the corresponding two-factor statement; see \cite{Ando1995}.

The multivariable conjecture in the form relevant here was recorded by Lin in the context of Hua determinant inequalities \cite{Lin2024}.  Pang subsequently disproved the conjecture in two complementary directions \cite{Pang2026}:
\begin{itemize}[leftmargin=2em]
\item for every $m\ge4$, failure already occurs for $2\times2$ real rank-one orthogonal projections, and the failure is at the largest eigenvalue;
\item for $m=3$, a $3\times3$ real positive semidefinite counterexample exists, with failure at the second eigenvalue.
\end{itemize}
Consequently the cell
\[
\boxed{m=3,\qquad n=2}
\]
is a natural unresolved boundary case: increasing either the number of factors or the matrix dimension is already known to allow failure.

\subsection{Target formulation for the project}
Throughout the project we put
\[
P=X^\alpha Y^\beta Z^\gamma,
\qquad
S=\alpha X+\beta Y+\gamma Z,
\]
and seek
\begin{equation}
 s_1(P)\le\lambda_1(S),
 \qquad
 s_2(P)\le\lambda_2(S),
 \label{eq:two-targets}
\end{equation}
for all $X,Y,Z\in\mathbb H_+^2$ and all positive weights $\alpha,\beta,\gamma$ summing to one.

\subsection{Current project-level assessment}
The work completed so far does not close \eqref{eq:two-targets} in full.  It does, however, remove several large boundary families and reduce any interior extremal counterexample to a small set of scalar parameters.  The equal-weight slice
\[
\alpha=\beta=\gamma=\frac13
\]
is substantially more rigid and is the recommended first target for completion.

\section{Closed boundary cases and general reductions}

\subsection{Commuting triples}
\begin{proposition}[Commuting case \statusR]
If $X,Y,Z$ commute, then \eqref{eq:MY} holds for both $j=1,2$.
\end{proposition}
\begin{proof}
The three matrices are simultaneously diagonalizable.  On each common eigenline the desired estimate becomes the scalar weighted Young inequality
\[
x^\alpha y^\beta z^\gamma\le \alpha x+\beta y+\gamma z.
\]
Ordering the two resulting scalar coordinates gives the required eigenvalue inequalities.
\end{proof}

\subsection{One scalar factor}
\begin{proposition}[Scalar-factor case \statusR]
If one of $X,Y,Z$ is a scalar matrix, then \eqref{eq:MY} holds in every dimension.
\end{proposition}
\begin{proof}
It suffices to take $Z=zI$.  Set
\[
\delta=\alpha+\beta,\qquad a=\frac{\alpha}{\delta},\qquad b=\frac{\beta}{\delta}.
\]
Then $a+b=1$ and
\[
X^\alpha Y^\beta=(X^\delta)^a(Y^\delta)^b.
\]
By Ando's two-factor Matrix Young inequality,
\[
s_j(X^\alpha Y^\beta)
\le \lambda_j\bigl(aX^\delta+bY^\delta\bigr).
\]
Because $0<\delta<1$, operator concavity of $t\mapsto t^\delta$ gives
\[
aX^\delta+bY^\delta\preceq (aX+bY)^\delta.
\]
Functional calculus and the scalar weighted Young inequality imply
\[
z^\gamma(aX+bY)^\delta
\preceq \delta(aX+bY)+\gamma zI
=\alpha X+\beta Y+\gamma zI.
\]
Combining the inequalities proves the claim.
\end{proof}

\subsection{Adjacent identical factors}
\begin{proposition}[Adjacent-equal case \statusR]
If $X=Y$ or $Y=Z$, then \eqref{eq:MY} reduces to the two-factor Ando inequality.
\end{proposition}
\begin{proof}
For example, if $X=Y$ then
\[
X^\alpha Y^\beta Z^\gamma=X^{\alpha+\beta}Z^\gamma,
\]
and $(\alpha+\beta)+\gamma=1$.
\end{proof}

\subsection{A determinant slack available in dimension two}
For the remainder of the manuscript let $n=2$ and assume positive definiteness when inverses or logarithms are used; semidefinite cases can be obtained by regularization.

Put
\[
P=X^\alpha Y^\beta Z^\gamma,
\qquad
S=\alpha X+\beta Y+\gamma Z.
\]
Define
\[
\rho=\sqrt{\det P},\qquad \Delta=\sqrt{\det S}.
\]
Since $\det^{1/2}$ is concave on $\mathbb H_+^2$,
\[
\Delta\ge \alpha\sqrt{\det X}+\beta\sqrt{\det Y}+\gamma\sqrt{\det Z}.
\]
Applying scalar weighted AM--GM to the three positive numbers $\sqrt{\det X}$, $\sqrt{\det Y}$ and $\sqrt{\det Z}$ yields
\begin{equation}
\boxed{\Delta\ge \rho.}
\label{eq:det-slack}
\end{equation}
This nonnegative determinant slack is one of the principal special features of the $2\times2$ problem.

\section{Dimension-two structure}

\subsection{Hyperbolic reformulation of the two eigenvalue inequalities}
\begin{proposition}[Exact two-by-two reformulation \statusR]
Let
\[
s_1(P)=\rho e^L,\qquad s_2(P)=\rho e^{-L},
\]
and
\[
\lambda_1(S)=\Delta e^h,\qquad \lambda_2(S)=\Delta e^{-h},
\]
where $L,h\ge0$.  Then the two desired inequalities in \eqref{eq:two-targets} are jointly equivalent to
\begin{equation}
\boxed{
|L-h|\le \log\frac{\Delta}{\rho}.
}
\label{eq:hyperbolic-slack}
\end{equation}
\end{proposition}
\begin{proof}
Since $s_1(P)s_2(P)=\det P=\rho^2$ and $\lambda_1(S)\lambda_2(S)=\det S=\Delta^2$, the displayed parametrizations are exact.  The inequalities
\[
\rho e^L\le \Delta e^h,
\qquad
\rho e^{-L}\le \Delta e^{-h}
\]
are respectively equivalent to
\[
L-h\le\log(\Delta/\rho),
\qquad
-(L-h)\le\log(\Delta/\rho).
\]
Combining them gives \eqref{eq:hyperbolic-slack}.  The right side is nonnegative by \eqref{eq:det-slack}.
\end{proof}

This formulation separates the problem into an anisotropy mismatch $|L-h|$ and a determinant slack $\log(\Delta/\rho)$.

\subsection{Scale-stationary normalization}
The following reduction is useful when studying a hypothetical interior maximizer of either ratio
\[
R_j(X,Y,Z)=\frac{s_j(X^\alpha Y^\beta Z^\gamma)}{\lambda_j(\alpha X+\beta Y+\gamma Z)}.
\]

\begin{lemma}[Stationary scaling condition \statusR]\label{lem:stationary-scaling}
Fix the shapes of three positive definite matrices $X_1,X_2,X_3$ and vary only positive scalar multipliers $c_i$.  Suppose $R_j(c_1X_1,c_2X_2,c_3X_3)$ has an interior stationary point and the relevant eigenvalue of
\[
S(c)=\sum_{i=1}^3\alpha_i c_iX_i
\]
is simple, with normalized eigenvector $v$.  Then
\begin{equation}
 c_i\langle v,X_iv\rangle=\lambda_j(S(c)),\qquad i=1,2,3.
\label{eq:stationary-rayleigh}
\end{equation}
\end{lemma}
\begin{proof}
The numerator is multiplied by $\prod_i c_i^{\alpha_i}$.  Hence the derivative of its logarithm with respect to $\log c_i$ is $\alpha_i$.  By Hellmann--Feynman,
\[
\frac{\partial}{\partial\log c_i}\log\lambda_j(S(c))
=\frac{\alpha_i c_i\langle v,X_iv\rangle}{\lambda_j(S(c))}.
\]
Stationarity of $\log R_j$ gives \eqref{eq:stationary-rayleigh}.
\end{proof}

After dividing all three scaled matrices by $\lambda_j(S(c))$ and choosing $v=e_1$, one may therefore write
\begin{equation}
B_i=
\begin{pmatrix}
1&a_i\\
\overline{a_i}&b_i
\end{pmatrix}>0,
\qquad
\sum_{i=1}^3\alpha_i a_i=0,
\label{eq:stationary-normal-form}
\end{equation}
and
\[
\sum_{i=1}^3\alpha_i B_i=\diag(1,s).
\]
For the $j=1$ problem one has $s\le1$; for the $j=2$ problem one has $s\ge1$.

\subsection{Bloch/Pauli parameterization}
Every Hermitian $2\times2$ matrix has a representation
\[
X=x_0I+x\,n_x\cdot\sigma,
\]
where $n_x\in S^2$ and $\sigma=(\sigma_1,\sigma_2,\sigma_3)$ denotes the Pauli matrices.  Thus a triple of fixed spectra depends geometrically only on the Gram variables
\[
u=n_x\cdot n_y,\qquad v=n_y\cdot n_z,\qquad w=n_x\cdot n_z,
\]
which satisfy
\begin{equation}
1+2uvw-u^2-v^2-w^2\ge0.
\label{eq:gram-region}
\end{equation}
The boundary of \eqref{eq:gram-region} is exactly the coplanar case and can be parameterized by two angles,
\[
u=\cos a,\qquad v=\cos b,\qquad w=\cos(a-b).
\]

\paragraph{Working reduction.}
\statusW{} A Hessian calculation carried out during the project indicates that, for fixed spectra and weights, an extremum of either violation ratio can be pushed to the Gram boundary, which would reduce the complex Hermitian problem to real symmetric triples.  This is a promising reduction, but the present manuscript does not yet promote it to theorem status: the Hessian argument should be independently rewritten and checked before publication.

\section{The equal-weight slice $\alpha=\beta=\gamma=1/3$}

The equal-weight problem is
\begin{equation}
 s_j(B_1^{1/3}B_2^{1/3}B_3^{1/3})
 \le
 \lambda_j\!\left(\frac{B_1+B_2+B_3}{3}\right),
 \qquad j=1,2.
 \label{eq:equal-weight}
\end{equation}
It is the most rigid nontrivial slice and should be closed before returning to arbitrary weights.

\subsection{Largest singular value: reduction to arithmetic mean $I$}
\begin{proposition}[Equal-weight $j=1$ normalization \statusR]
Assume a $j=1$ counterexample occurs at an interior scale-stationary point.  After the normalization of Lemma~\ref{lem:stationary-scaling}, one may further reduce to
\begin{equation}
\boxed{B_1+B_2+B_3=3I.}
\label{eq:sum3I}
\end{equation}
\end{proposition}
\begin{proof}
At a scale-stationary point, choose the top eigenvector of the arithmetic mean as $e_1$.  Then
\[
B_i=\begin{pmatrix}1&q_i\\ \overline q_i&b_i\end{pmatrix},
\qquad
\frac{B_1+B_2+B_3}{3}=\diag(1,s),
\qquad s\le1.
\]
If $s<1$, define
\[
B_1'=B_1+3(1-s)E_{22}.
\]
Then $(B_1'+B_2+B_3)/3=I$ and $B_1'\succeq B_1$.  Since $t\mapsto t^{2/3}$ is operator monotone,
\[
B_1'^{\,2/3}\succeq B_1^{2/3}.
\]
For
\[
P=B_1^{1/3}B_2^{1/3}B_3^{1/3},
\]
we have
\[
P^*P=B_3^{1/3}B_2^{1/3}B_1^{2/3}B_2^{1/3}B_3^{1/3}.
\]
Replacing $B_1$ by $B_1'$ therefore cannot decrease $s_1(P)$.  Hence a normalized counterexample with $s<1$ produces one with arithmetic mean exactly $I$.
\end{proof}

Under \eqref{eq:sum3I}, write
\[
B_i=\begin{pmatrix}1&q_i\\ \overline q_i&|q_i|^2+r_i\end{pmatrix},
\qquad r_i=\det B_i>0.
\]
The constraints become
\begin{equation}
q_1+q_2+q_3=0,
\qquad
|q_1|^2+|q_2|^2+|q_3|^2+r_1+r_2+r_3=3.
\label{eq:equal-j1-constraints}
\end{equation}
Thus even before any real-symmetric reduction, the scale-normalized equal-weight $j=1$ problem is a finite scalar problem.

\subsection{A candidate dimension-free strengthening}
The following inequality emerged naturally during the project:
\begin{problem}[Candidate three-factor norm inequality \statusW]
For positive semidefinite matrices $A,B,C$, prove or disprove
\begin{equation}
\boxed{
\|ABC\|\le \frac13\|A^3+B^3+C^3\|.
}
\label{eq:candidate-abc}
\end{equation}
\end{problem}
If true, \eqref{eq:candidate-abc} would settle the equal-weight $j=1$ inequality in every dimension, not merely in dimension two.

A valid Bhatia--Davis matrix AGM estimate gives
\begin{equation}
2\|ABC\|\le \|A^2B+BC^2\|.
\label{eq:bd-bridge}
\end{equation}
The missing step is a sharp upper estimate connecting the right side of \eqref{eq:bd-bridge} to $A^3+B^3+C^3$.  Several attempted bridges were explored, but no completely audited proof of the constant $2/3$ has yet been retained.  Therefore \eqref{eq:candidate-abc} is deliberately \emph{not} promoted to theorem status in this manuscript.

\subsection{Smallest singular value: a sharp trace-sign reduction}
Now consider a scale-stationary candidate counterexample for $j=2$.  Normalize as in \eqref{eq:stationary-normal-form} so that
\begin{equation}
B_i=\begin{pmatrix}1&q_i\\ \overline q_i&|q_i|^2+r_i\end{pmatrix}>0,
\qquad q_1+q_2+q_3=0,
\label{eq:j2-normal}
\end{equation}
and
\[
\frac{B_1+B_2+B_3}{3}=\diag(1,s),\qquad s\ge1.
\]
Set
\[
P=B_1^{1/3}B_2^{1/3}B_3^{1/3},
\qquad
d=\det P=(r_1r_2r_3)^{1/3},
\qquad
\tau=\Tr(P^*P).
\]

\begin{proposition}[Safe determinant region \statusR]
If $d\le1$, then $s_2(P)\le1$ automatically.
\end{proposition}
\begin{proof}
Since $s_1(P)s_2(P)=d$ and $s_1(P)\ge s_2(P)$,
\[
s_2(P)\le \sqrt d\le1.
\]
\end{proof}

Thus a normalized $j=2$ counterexample must satisfy $d>1$.  Moreover Minkowski's determinant inequality gives
\[
\sqrt{s}=\sqrt{\det\!\left(\frac{B_1+B_2+B_3}{3}\right)}
\ge \frac{\sqrt{r_1}+\sqrt{r_2}+\sqrt{r_3}}{3}
\ge (r_1r_2r_3)^{1/6}=\sqrt d,
\]
so
\begin{equation}
1<d\le s.
\label{eq:d-region}
\end{equation}

\begin{proposition}[Trace-sign equivalence for the remaining $j=2$ region \statusR]
In the region $d>1$, the desired inequality $s_2(P)\le1$ is equivalent to
\begin{equation}
\boxed{
\tau\ge1+d^2.
}
\label{eq:J2}
\end{equation}
\end{proposition}
\begin{proof}
Let the eigenvalues of $P^*P$ be $x=s_1(P)^2\ge y=s_2(P)^2$, so $xy=d^2>1$.  Hence $x>1$.  Now
\[
\det(I-P^*P)=(1-x)(1-y)=1-\tau+d^2.
\]
Because $1-x<0$, the condition $y\le1$ is equivalent to $\det(I-P^*P)\le0$, which is exactly \eqref{eq:J2}.
\end{proof}

Equation \eqref{eq:J2}, under \eqref{eq:j2-normal} and \eqref{eq:d-region}, is the principal rigorous bottleneck currently remaining in the equal-weight $j=2$ problem.

\subsection{Numerical status}
\statusN{} Extensive random and reduced-parameter searches performed during the project have not found a violation of either equal-weight inequality in dimension two.  The experiments also suggest that equality is concentrated on commuting or degenerate boundary configurations.  These observations are guidance only and are not used as proofs.

\section{General weights: what remains after the reductions}

For arbitrary positive $\alpha,\beta,\gamma$ with sum one, the project has not yet found a single identity as rigid as the cubic equal-weight structure.  Nevertheless, the general problem has already been reduced in several useful ways.

\subsection{Interior extremizers have a common Rayleigh normalization}
By the scale-stationary condition, a hypothetical interior extremizer can be normalized so that
\[
B_i=\begin{pmatrix}1&a_i\\ \overline a_i&b_i\end{pmatrix},
\qquad
\alpha a_1+\beta a_2+\gamma a_3=0,
\]
and
\[
\alpha B_1+\beta B_2+\gamma B_3=\diag(1,s).
\]
For the top-eigenvalue problem $s\le1$, whereas for the bottom-eigenvalue problem $s\ge1$.  Thus any genuinely new obstruction must already survive a shared Rayleigh-quotient normalization.

\subsection{Two-dimensional determinant coupling}
The identity
\[
s_1(P)s_2(P)=\det P
\]
means that the $j=1$ and $j=2$ inequalities are not independent in dimension two.  The exact formulation \eqref{eq:hyperbolic-slack} says that both are controlled by a single anisotropy mismatch.  A successful estimate for $L-h$ that uses the determinant slack \eqref{eq:det-slack} would close both indices at once.

\subsection{Rayleigh-envelope candidates}
During the project the following stronger inequalities were identified as natural candidates:
\begin{align}
 s_1(X^\alpha Y^\beta Z^\gamma)
 &\stackrel{?}{\le}
 \max_{\|v\|=1}
 (v^*Xv)^\alpha(v^*Yv)^\beta(v^*Zv)^\gamma,
 \label{eq:Eplus}\\
 s_2(X^\alpha Y^\beta Z^\gamma)
 &\stackrel{?}{\le}
 \min_{\|v\|=1}
 (v^*Xv)^\alpha(v^*Yv)^\beta(v^*Zv)^\gamma.
 \label{eq:Eminus}
\end{align}
If both held, scalar weighted Young applied pointwise in $v$ would prove the Matrix Young inequalities immediately.  However, these envelope statements are stronger than the original target and are currently \statusW{} rather than established lemmas.  They should not be used as hidden assumptions.

\subsection{Rank-one boundary}
\statusW{} The project derived a geometric argument suggesting that three rank-one factors satisfy the $m=3,n=2$ inequality, in contrast with Pang's $m\ge4$ rank-one counterexample mechanism.  The argument passes through the projective geometry of three lines in $\mathbb C^2$.  Because the covering-radius step has not yet received a separate proof audit, this result is recorded as a working lemma rather than as a closed theorem in the present source.

\subsection{Relation to Hua determinant inequalities}
Lin's 2024 paper develops Young-type eigenvalue inequalities alongside extensions of the Hua determinant inequality \cite{Lin2024}.  Its discussion of generalized H\"older eigenvalue inequalities for contractive matrix products is structurally close to the determinant-sign expression
\[
\det(I-P^*P)=1-\Tr(P^*P)+\det(P)^2
\]
that appears in the normalized $2\times2$ $j=2$ problem.  A serious next literature step is therefore to test whether an appropriate contractive rescaling converts the remaining trace-sign inequality into a direct consequence of the Hua/H\"older machinery, rather than expanding matrix cube roots by hand.

\subsection{Current general-weight bottlenecks}
At present the general-weight problem has two main unresolved parts:
\begin{enumerate}[leftmargin=2.2em]
\item prove the top singular-value inequality for genuinely noncommuting, full-rank $2\times2$ triples away from the scalar and commuting boundaries;
\item transfer that control, together with determinant slack, to the bottom singular value, or derive a direct normalized trace-sign inequality analogous to \eqref{eq:J2}.
\end{enumerate}
The equal-weight case should be treated as the model problem from which the correct general-weight scalar inequality is to be inferred.

\section{Audit: routes that must not be silently reused}

This section records approaches explored during the project that are either false or not yet justified.  Keeping them in the source is important because several of them look deceptively natural.

\subsection{A false Loewner bridge}
Let $\delta=\beta+\gamma$ and consider
\[
H_{\beta,\gamma}(Y,Z)
=\left(Y^\beta Z^{2\gamma}Y^\beta\right)^{1/(2\delta)}.
\]
A tempting strategy was to combine Ando's two-factor inequality with
\begin{equation}
H_{\beta,\gamma}(Y,Z)
\stackrel{?}{\preceq}
\frac{\beta Y+\gamma Z}{\beta+\gamma}.
\label{eq:false-loewner}
\end{equation}
\statusF{} This Loewner inequality is false even for $2\times2$ positive definite matrices; numerical counterexamples were found during the audit.  Hence no proof should invoke \eqref{eq:false-loewner}.  An exact rational/algebraic counterexample should eventually be stored in an appendix for reproducibility.

\subsection{``The angular maximizer is commuting''}
Another attempted shortcut was to show that, for fixed spectra, the violation ratio is maximized at commuting orientations.  \statusF{} Reduced two-angle optimization finds noncommuting local and sometimes global maximizers for fixed spectra, although the observed maximal ratio remains below one.  Thus ``the worst point is commuting'' is too strong and cannot serve as a proof strategy.

\subsection{Premature use of a numerical-radius bridge}
The candidate dimension-free estimate \eqref{eq:candidate-abc} suggested trying to strengthen Bhatia--Davis from
\[
2\|ABC\|\le\|A^2B+BC^2\|
\]
to a numerical-radius estimate.  \statusW{} No fully audited argument for the required numerical-radius bridge has been retained.  In particular, operator-norm and numerical-radius bounds must not be interchanged without a separate theorem.

\subsection{Withdrawn rectangular AGM manipulation}
\statusF{} An earlier scratch derivation used a rectangular arithmetic--geometric mean inequality with the Gram factors on the wrong side.  That step was explicitly withdrawn during the audit.  The only Bhatia--Davis form used in the retained manuscript is the standard square-matrix estimate leading to \eqref{eq:bd-bridge}; see \cite{BhatiaDavis1993}.

\subsection{Numerics are not a classification theorem}
\statusN{} Large random searches, ill-conditioned tests, and reduced two-angle global optimization have so far found no $m=3,n=2$ counterexample.  This strongly motivates the positive conjecture but does not rule out a thin algebraic counterexample family.  All formal claims in this manuscript are therefore independent of these computations.

\section{Work map and publication-level gaps}

The project should now be advanced in the following order.

\subsection*{Priority A: close the equal-weight $j=2$ trace-sign inequality}
The sharp remaining normalized statement is \eqref{eq:J2}:
\[
\Tr(P^*P)\ge 1+d^2,
\qquad
P=B_1^{1/3}B_2^{1/3}B_3^{1/3},
\qquad
1<d\le s,
\]
under \eqref{eq:j2-normal}.  There are two serious routes:
\begin{enumerate}[leftmargin=2.2em]
\item \textbf{Hua/H\"older route.}  Find a contractive rescaling that converts $\det(I-P^*P)$ into an instance of the generalized H\"older eigenvalue inequalities discussed around Hua determinant inequalities; compare \cite{Lin2024}.
\item \textbf{Explicit $2\times2$ route.}  Use Cayley--Hamilton to write
\[
B_i^{1/3}=a_iI+b_iB_i
\]
and expand $\Tr(P^*P)-1-d^2$ under $q_1+q_2+q_3=0$.  Search for a decomposition into commutator squares and manifestly nonnegative scalar terms.
\end{enumerate}

\subsection*{Priority B: independently audit the equal-weight $j=1$ candidate}
Prove or disprove \eqref{eq:candidate-abc}.  If it is true, identify the shortest legitimate bridge from Bhatia--Davis or from an Araki--Lieb--Thirring/three-factor AGM inequality.  Israel--Krahmer--Ward prove a noncommutative AGM for averaged products of up to three matrices \cite{IsraelKrahmerWard2016}; their result is related but is not identical to the fixed ordered-product inequality needed here.

If \eqref{eq:candidate-abc} fails in higher dimension but survives for $2\times2$, the failure mechanism itself may explain why $(m,n)=(3,2)$ is exceptional.

\subsection*{Priority C: complete the real-symmetric reduction audit}
Rewrite the Pauli/Gram Hessian argument from scratch.  Either prove that a maximal violation ratio can always be pushed to the boundary
\[
1+2uvw-u^2-v^2-w^2=0,
\]
or produce a counterexample to that reduction.  Once certified, the remaining angular geometry has only two real variables.

\subsection*{Priority D: return to arbitrary weights}
After the equal-weight case is settled, use the successful proof to determine which features survive when $1/3$ is replaced by $(\alpha,\beta,\gamma)$.  The likely final mechanisms are:
\begin{itemize}[leftmargin=2em]
\item determinant slack plus anisotropy control via \eqref{eq:hyperbolic-slack};
\item a weighted analogue of the equal-weight trace-sign inequality;
\item or a weighted two-angle scalar inequality after real-symmetric reduction.
\end{itemize}

\subsection{What would constitute a complete paper?}
A publication-ready result need not necessarily prove the entire general-weight conjecture.  The following outcomes are already mathematically coherent paper targets:
\begin{enumerate}[leftmargin=2.2em]
\item a proof of the full $(m,n)=(3,2)$ cell, giving the sharp positive boundary between Pang's two counterexample regimes;
\item a complete equal-weight theorem together with equality classification and a structural $2\times2$ proof;
\item a counterexample in the remaining cell, together with the minimal parameter mechanism;
\item or a new dimension-free three-factor norm inequality such as \eqref{eq:candidate-abc}, if independently proved, plus its relation to Matrix Young.
\end{enumerate}

\subsection{Current status table}
\begin{center}
\begin{tabularx}{\textwidth}{@{}l l X@{}}
\toprule
Component & Status & Main remaining issue\\
\midrule
Literature boundary & \statusL{} & none\\
Commuting / scalar / adjacent-equal & \statusR{} & equality audit only\\
Determinant-hyperbolic reformulation & \statusR{} & none\\
Scale-stationary normalization & \statusR{} & boundary/non-simple cases to package\\
Complex $\to$ real-symmetric extremizer & \statusW{} & independent Hessian audit\\
Equal weight, $j=1$ & \statusW{} & prove/disprove \eqref{eq:candidate-abc}\\
Equal weight, $j=2$ & \statusR{} reduction & prove trace-sign \eqref{eq:J2}\\
General weights, $j=1$ & open & main global inequality\\
General weights, $j=2$ & open & determinant transfer / weighted trace sign\\
Numerical search & \statusN{} & no counterexample found\\
\bottomrule
\end{tabularx}
\end{center}


\bibliographystyle{alpha}
\bibliography{references}

\end{document}
